\documentclass[11pt]{article}
\usepackage[T1]{fontenc}
\usepackage[utf8]{inputenc}
\usepackage{lmodern,amsmath,amssymb,amsthm,booktabs,array}
\usepackage[margin=1in]{geometry}
\usepackage[colorlinks=true,urlcolor=blue,citecolor=blue,linkcolor=blue]{hyperref}
\newtheorem{theorem}{Theorem}
\newtheorem{lemma}[theorem]{Lemma}
\newtheorem{corollary}[theorem]{Corollary}
\title{An elementary consequence of Baishanski and Neuwirth\\
answering Hanna's coefficient root-limit question}
\author{Haobo Ma\\\small ChronoAI Pte Ltd, Singapore
\and Wenlin Zhang\\\small National University of Singapore\\
\small ChronoAI Pte Ltd, Singapore}
\date{November 14, 2026}
\begin{document}
\maketitle
\begin{abstract}
As an elementary consequence of Baishanski's coefficient-norm limit
(1999, \S28.1, p.~523) and Neuwirth's trinomial maximum
(2007, Remark~6.4), we answer Paul D. Hanna's question in OEIS A195417,
first comment (history revision~2, September 17, 2011),
for the cubic polynomial $1+x-x^3$.
If $a_n$ is the sum of the absolute values of the coefficients of
$(1+x-x^3)^n$, then
\[
 \lim_{n\to\infty}a_n^{1/n}
 =\sqrt{\frac{97+26\sqrt{13}}{27}}
 =2.6579324543577637\ldots.
\]
A unit-circle coefficient estimate gives explicit two-sided bounds for
every $n$, and a polynomial identity gives the maximum exactly.
The same estimate gives a self-contained proof of the A084611
root limit $\sqrt5$, announced by Kot\v{e}\v{s}ovec in 2013,
and gives the unposed root limits
$9/\sqrt8$ and $\sqrt{52/3}$ for A084613 and A084615.
\end{abstract}

\section{The question and its answer}
For a polynomial $Q$, write $[x^k]Q$ for its coefficient of $x^k$.
Define
\[
 P(x)=1+x-x^3,\qquad
 a_n=\sum_{k=0}^{3n}\bigl|[x^k]P(x)^n\bigr|
 \quad(n=0,1,2,\ldots).
\]
In the first comment of A195417, Hanna asks, ``What is the limit
a(n)\^{}(1/n)?''~\cite{hanna195417}. The answer is
$R=\sqrt{(97+26\sqrt{13})/27}$. The question and the coefficient definition
already appear in the September 17, 2011 entry history, revision~2.
The cancellation of signed coefficients prevents evaluation at $x=1$
from determining their absolute sum. The relevant maximum is instead
on the complex unit circle.

\begin{theorem}\label{thm:main}
Put $R=\sqrt{(97+26\sqrt{13})/27}$. For every integer $n\geq0$,
\begin{equation}\label{eq:bounds}
 R^n\leq a_n\leq\sqrt{3n+1}\,R^n.
\end{equation}
Consequently, $\lim_{n\to\infty}a_n^{1/n}=R$.
\end{theorem}

The cross-referenced entry A084611 concerns $1+x-x^2$ and conjectures
the root limit $\sqrt5$~\cite{hanna84611}. Kot\v{e}\v{s}ovec's note dated
July 26, 2013 announces that limit and a Legendre-polynomial
treatment~\cite{kotesovec}.
Corollary~\ref{cor:quadratic} supplies an alternative proof within the
same coefficient estimate.

\paragraph{Earlier results and the present argument.}
The coefficient-norm root limit is classical: Girard discusses it in
1971~\cite[\S1, pp.~357--358]{girard}, and Baishanski explicitly states
the same limit for absolutely convergent power series
in 1999~\cite[\S28.1, p.~523]{baishanski}.
Neuwirth's 2007 formulas give both the cubic and quadratic circle
maxima used here~\cite[Remarks~6.3--6.4, pp.~13--14]{neuwirth}.
Thus those earlier results already imply the displayed answer by
specialization. The contribution here is a self-contained elementary
derivation with explicit bounds for every $n$ and short algebraic identities.
Calegari explicitly printed the same squared cubic maximum in
2014 for the equivalent Laurent polynomial $z^2+z-z^{-1}$~\cite{calegari}.
The later convolution-power theorems of Diaconis--Saloff-Coste
(May 2012) and Randles--Saloff-Coste (December 2012) also imply these
root limits~\cite{diaconis,randles}; a stronger bound from the former
is described below.

\section{A coefficient estimate}
The following finite-degree estimate is a direct Parseval--Cauchy--Schwarz
proof of the polynomial case of Baishanski's stated limit~\cite{baishanski}.
\begin{lemma}\label{lem:general}
Let $Q\in\mathbb C[x]$ be nonzero of degree $d$, and define
\[
 M=\max_{|z|=1}|Q(z)|,\qquad
 L_n=\sum_{k=0}^{dn}|[x^k]Q(x)^n|.
\]
Then $M>0$, and for every integer $n\geq0$,
\begin{equation}\label{eq:general}
 M^n\leq L_n\leq\sqrt{dn+1}\,M^n.
\end{equation}
In particular $\lim_{n\to\infty}L_n^{1/n}=M$.
\end{lemma}
\begin{proof}
Continuity on the unit circle gives a maximizing point $z_0$.
Since a nonzero polynomial has only finitely many zeros, $M>0$.
Write $Q(x)^n=\sum_{k=0}^{dn}c_kx^k$.
The triangle inequality at $z_0$ gives
$M^n=|\sum_kc_kz_0^k|\leq\sum_k|c_k|=L_n$.
For the upper bound, expand the finite product under the integral.
The identity
\[
 \frac1{2\pi}\int_0^{2\pi}e^{\mathrm i(k-\ell)t}\,dt
 =\begin{cases}1&k=\ell,\\0&k\ne\ell\end{cases}
\]
gives
\[
 \sum_{k=0}^{dn}|c_k|^2
 =\frac1{2\pi}\int_0^{2\pi}|Q(e^{\mathrm it})|^{2n}\,dt
 \leq M^{2n}.
\]
Cauchy--Schwarz now yields
$L_n\leq\sqrt{dn+1}(\sum_k|c_k|^2)^{1/2}$, proving
\eqref{eq:general}, also for $n=0$.
For $n\geq1$, take $n$th roots and use
$(dn+1)^{1/(2n)}\to1$. This proves the limit.
\end{proof}

\section{The cubic maximum}
For comparison with Neuwirth's Remark~6.4, put
$T(x)=e^{-2\mathrm ix}+e^{\mathrm i\pi/3}+e^{\mathrm ix}$.
Then
\[
 P(e^{\mathrm it})=e^{\mathrm i(t-\pi/3)}T(\pi/3-t).
\]
His formula with $r_1=r_2=r_3=1$ gives
$\max|T|^2=11/3+2((13/9)^{3/2}-1/27)=(97+26\sqrt{13})/27$.
The factorization below supplies a direct algebraic proof of this
specialization. Calegari's equivalent expression follows from
$z^2+z-z^{-1}=z^2P(z^{-1})$ on the unit circle~\cite{calegari}.
\begin{proof}[Proof of Theorem~\ref{thm:main}]
For $|z|=1$ put $u=\operatorname{Re}z\in[-1,1]$.
Expanding $P(z)\overline{P(z)}$ and using
$\operatorname{Re}(z^2)=2u^2-1$ and
$\operatorname{Re}(z^3)=4u^3-3u$ gives
\begin{equation}\label{eq:cubic}
 |1+z-z^3|^2=5+8u-4u^2-8u^3.
\end{equation}
Set $s=\sqrt{13}$. The exact identity
\begin{equation}\label{eq:gap}
 \frac{97+26s}{27}-(5+8u-4u^2-8u^3)
 =\frac{(6u+1-s)^2(6u+1+2s)}{27}
\end{equation}
holds because $s^2=13$.
Its right side is nonnegative on $[-1,1]$: its last factor
is at least $2s-5>0$. Equality holds at
$u_0=(s-1)/6\in[-1,1]$.
The point $z_0=u_0+\mathrm i\sqrt{1-u_0^2}$ lies on the unit circle
and attains equality. Thus $\max_{|z|=1}|P(z)|=R$.
Apply Lemma~\ref{lem:general} with $d=3$ to obtain
\eqref{eq:bounds} and the asserted limit.
\end{proof}

The estimate also gives a quantitative convergence bound:
for $n\geq1$,
\[
 0\leq\log(a_n^{1/n})-\log R\leq\frac{\log(3n+1)}{2n}.
\]
This follows by taking logarithms in \eqref{eq:bounds}.

\section{The cross-referenced quadratic family}
\begin{corollary}\label{cor:quadratic}
For a real $m\geq1$, put
\[
 B_n(m)=\sum_{k=0}^{2n}|[x^k](1+x-mx^2)^n|,\qquad
 M_m=(m+1)\sqrt{1+\frac1{4m}}.
\]
Then, for every integer $n\geq0$,
\[
 M_m^n\leq B_n(m)\leq\sqrt{2n+1}\,M_m^n,
 \qquad\lim_{n\to\infty}B_n(m)^{1/n}=M_m.
\]
In particular the limits for A084611, A084613 and A084615 are, respectively,
$\sqrt5$, $9/\sqrt8$ and $\sqrt{52/3}$.
\end{corollary}
\begin{proof}
Writing $u=\operatorname{Re}z$ for $|z|=1$ gives
\[
 |1+z-mz^2|^2=m^2+2m+2+2(1-m)u-4mu^2
 =M_m^2-\frac{(4mu+m-1)^2}{4m}.
\]
Equality holds at $u=(1-m)/(4m)\in[-1/4,0]$.
Lemma~\ref{lem:general} with $d=2$ now gives the bounds and limit.
\end{proof}

The A084613 and A084615 entries define their sequences but pose no
root-limit conjecture~\cite{hanna-family}; the corresponding limits above
are consequences of those definitions, not answers to asserted conjectures.
The maximum $M_m$ also follows from Neuwirth's Remark~6.3 with
$(r_1,r_2,r_3)=(1,1,m)$: for $x=t-\pi/2$,
$e^{-\mathrm ix}+\mathrm i+m e^{\mathrm ix}
=\mathrm i e^{-\mathrm it}(1+e^{\mathrm it}-m e^{2\mathrm it})$.
The required inequality $|1-1/m|<4$ holds for $m\geq1$.

\paragraph{A stronger non-explicit bound from earlier stability theory.}
For each fixed polynomial $Q=P$ or $Q(x)=1+x-mx^2$, $m\geq1$,
Diaconis--Saloff-Coste's Theorem~2.1 gives a constant $C_Q$ such that
$\sum_k|[x^k]Q^n|\leq C_Q(\max_{|z|=1}|Q(z)|)^n$~\cite{diaconis}.
Here is the hypothesis check. The squared moduli have exactly two
circle maxima. At either cubic maximum their second derivative with
respect to $t$ is $-8\sqrt{13}(1-u_0^2)<0$; for the quadratic it is
$-8m(1-u_m^2)<0$, where $u_m=(1-m)/(4m)$.
Normalize the finite coefficient sequence by its circle maximum.
Near either maximizing angle $t_0$, a local analytic logarithm of its
Fourier transform has the expansion
$\mathrm i\alpha(t-t_0)-\beta(t-t_0)^2+O((t-t_0)^3)$
after division by its value at $t_0$, with $\alpha\in\mathbb R$
and $\operatorname{Re}\beta>0$. This is condition~2 of that theorem
with $\nu=2$, so the normalized convolution powers have bounded
coefficient norm. The square-root bounds above instead give explicit
constants for every $n$. In particular, Example~2.4 and the second
remark after Proposition~2.5 of that paper already yield the quadratic
root limit for $(a_0,a_+,a_-)=(1,-1,1)$, before the 2013 note.

\begin{corollary}[Row identities]\label{cor:rows}
For $T_{b,c}(n,k)=[x^k](1+bx+cx^2)^n$,
\begin{align*}
 T_{b,c}(n,k)&=\sum_{\substack{j\geq0\\0\leq k-2j\leq n-j}}
   \binom nj\binom{n-j}{k-2j}b^{k-2j}c^j,\\
 \sum_k T_{b,c}(n,k)&=(1+b+c)^n,\qquad
 \sum_k(-1)^k T_{b,c}(n,k)=(1-b+c)^n,\\
 T_{b,c}(n,2n)&=c^n,\qquad
 T_{b,c}(n,2n-1)=nbc^{n-1}\quad(n\geq1).
\end{align*}
For $b=1$, $c=-m$, the ordinary and alternating row sums are
$(2-m)^n$ and $(-m)^n$. For $b,c\geq0$ the absolute row sum is exactly
$(1+b+c)^n$. Powers with exponent zero have value one, including $0^0$.
\end{corollary}
\begin{proof}
Choose $j$ factors contributing $cx^2$ and $k-2j$ factors contributing $bx$.
Evaluation at $x=1,-1$ gives the row sums. The last two coefficients
come from taking $cx^2$ in all factors, or $bx$ once and $cx^2$ otherwise.
Nonnegative coefficients turn the ordinary sum into the absolute sum.
\end{proof}

Thus the nonnegative triangles A084600, A084602, A084604, A084606 and
A084608 have row sums $4^n,5^n,6^n,5^n,6^n$, respectively.
The signed triangles A084610, A084612 and A084614 give the three
absolute-sum sequences in Corollary~\ref{cor:quadratic}. These row
identities also give elementary proofs of their displayed sum and edge
formulas~\cite{hanna-triangles}.

\begin{corollary}[A diagonal sum in A084610]\label{cor:diagonal}
For $T(n,k)=[x^k](1+x-x^2)^n$,
\[
 \sum_{k=0}^nT(n-k,k)=\left\lfloor\frac{n+2}{2}\right\rfloor.
\]
\end{corollary}
\begin{proof}
The generating function of the left side is
$\sum_{r\geq0}t^r(1+t-t^2)^r
=1/((1-t)^2(1+t))=1/((1-t)(1-t^2))$.
The coefficient of $t^n$ counts the choices $0\leq j\leq\lfloor n/2\rfloor$.
\end{proof}

\begin{corollary}[A necessary ratio limit]\label{cor:ratio}
If the successive ratios of $a_n$, or of any $B_n(m)$ above, converge
to a positive number, that number equals the corresponding root limit.
\end{corollary}
\begin{proof}
For a positive sequence $u_n$, write
$\log u_n=\log u_0+\sum_{j=0}^{n-1}\log(u_{j+1}/u_j)$.
If the ratios tend to $\ell>0$, Ces\`aro averaging gives
$\log u_n/n\to\log\ell$.
\end{proof}
The root bounds establish this necessary value; the two source statements
about nonexistence of successive-ratio limits require a separate argument.

\section{Source coverage}
\begin{center}
\small
\begin{tabular}{@{}p{.31\linewidth}p{.61\linewidth}@{}}
\toprule
Source item & Mathematical coverage\\
\midrule
A195417, first comment & Answered: yes, the limit is $R$
(Theorem~\ref{thm:main}); consequence of Baishanski (1999) and
Neuwirth (2007), with an elementary proof here.\\[3pt]
A195417, second comment & The assertion about successive ratios and
attractors is not covered by the root-limit theorem.\\[3pt]
A084611, root-limit conjecture & Proved again in Corollary~\ref{cor:quadratic};
earlier consequence of Baishanski and Neuwirth; implied also by
Diaconis--Saloff-Coste (2012); announced by Kot\v{e}\v{s}ovec (2013).\\[3pt]
A084611, successive ratios & Not covered by these bounds.\\
A084613 and A084615, definitions & Not posed; the root limits are proved
in Corollary~\ref{cor:quadratic}, as specializations of earlier
norm-limit and maximum formulas.\\[3pt]
A084610, A084612, A084614, sum and edge formulas & Proved in
Corollary~\ref{cor:rows}; the A084610 diagonal sum is proved in
Corollary~\ref{cor:diagonal}.\\[3pt]
A084600, A084602, A084604, A084606, A084608, row sums & Proved in
Corollary~\ref{cor:rows}. Other identities and interpretations are not covered.\\[3pt]
A002426, A084601, A084603, A084605, A084609, A006139
(A084607 duplicate) & Central-coefficient asymptotics, congruences,
transforms and other source statements: not covered.\\
\bottomrule
\end{tabular}
\end{center}
The conditional ratio restriction in Corollary~\ref{cor:ratio} is not
asserted in either source; it neither proves nor refutes their
nonexistence statements.

\section{Verification scope}
The proofs above establish the coefficient estimate, the sharper square-root
prefactors, and both root limits. The Lean declaration
\texttt{D5.S3.Analytic.Polynomial.A195417RootLimit.result}
in \href{https://github.com/the-omega-institute/trureturing/blob/6bdbb68454659fcccffc7607b5fc292b2ff2fcdf/D5/S3/Analytic/Polynomial/A195417RootLimit.lean}{\texttt{A195417RootLimit.lean}}~\cite{lean}
states the cubic root limit for the
same integer coefficient sum, including the index $n=0$ in its defining
sequence. Its supporting upper estimate has prefactor $3n+1$;
the square-root prefactor, quadratic family and row identities have
written proofs above. Its definitions \texttt{p} and \texttt{a}
are the integer polynomial $1+X-X^3$ and
$\sum_{i=0}^{3n}|[X^i]p^n|$; the limit theorem uses positive indices
eventually, so its formal value at index zero has no effect on the limit.
The accompanying Python program checks exact polynomial identities,
initial coefficients, and finite instances of the inequalities; those
finite checks are distinct from the proofs for all $n$.
For the cubic and each of the three integer quadratic specializations,
the checked inequalities cover $0\leq n\leq1000$, and the independent
multinomial coefficient comparisons cover $0\leq n\leq29$.

\begin{thebibliography}{99}
\bibitem{girard}
D.~M. Girard, The asymptotic behavior of norms of powers of absolutely
convergent Fourier series, \emph{Pacific J. Math.} \textbf{37} (1971),
357--381, \S1.
\href{https://doi.org/10.2140/pjm.1971.37.357}{doi:10.2140/pjm.1971.37.357}.

\bibitem{baishanski}
B.~Baishanski, Norms of powers and a central limit theorem for
complex-valued probabilities, in \emph{Analysis of Divergence},
Birkh\"auser Boston, 1999, pp.~523--543, Chapter~28, \S28.1.
\href{https://doi.org/10.1007/978-1-4612-2236-1_29}{doi:10.1007/978-1-4612-2236-1\_29}.

\bibitem{neuwirth}
S.~Neuwirth, \emph{The maximum modulus of a trigonometric trinomial},
March 8, 2007, Remarks~6.3--6.4, pp.~13--14.
\href{https://arxiv.org/abs/math/0703236v1}{arXiv:math/0703236v1}.
Published in \emph{J. Anal. Math.} \textbf{104} (2008), 371--396,
\href{https://doi.org/10.1007/s11854-008-0028-2}{doi:10.1007/s11854-008-0028-2}.

\bibitem{diaconis}
P.~Diaconis and L.~Saloff-Coste, \emph{Convolution powers of complex
functions on $\mathbb Z$}, May 29, 2012, Theorem~2.1, Example~2.4,
and Proposition~2.5 with its following remarks.
\href{https://arxiv.org/abs/1205.6490v1}{arXiv:1205.6490v1}.

\bibitem{randles}
E.~Randles and L.~Saloff-Coste, \emph{On the convolution powers of
complex functions on $\mathbb Z$}, December 19, 2012, Theorem~1, p.~4.
\href{https://arxiv.org/abs/1212.4700v1}{arXiv:1212.4700v1}.

\bibitem{calegari}
F.~Calegari, \emph{The abelian house is not closed}, August 3, 2014;
collected in \emph{Persiflage: Math Blog Posts}, \S52.12, pp.~104--105.
\href{https://math.uchicago.edu/~fcale/papers/blog.pdf}{Author's collected posts}.

\bibitem{hanna195417}
P.~D. Hanna, A195417, \emph{The On-Line Encyclopedia of Integer Sequences},
September 17, 2011; comments 1--2, examples, and
\href{https://oeis.org/history?seq=A195417}{history, revision 2}.
\url{https://oeis.org/A195417}.

\bibitem{hanna84611}
P.~D. Hanna, A084611, \emph{The On-Line Encyclopedia of Integer Sequences},
June 1, 2003; root-limit comment added July 12, 2013
(history revision~7); Kot\v{e}\v{s}ovec link added July 28, 2013
(revision~9).
\url{https://oeis.org/A084611}.

\bibitem{kotesovec}
V.~Kot\v{e}\v{s}ovec, \emph{Asymptotic of sequence A084611},
July 26, 2013, pp.~1--3.
\url{http://www.kotesovec.cz/math_articles/kotesovec_asymptotic_of_sequence_A084611.pdf}.

\bibitem{hanna-family}
P.~D. Hanna, \emph{OEIS} A084613 and A084615, definitions,
June 1, 2003; versions 17 and 13, March 2023.
\url{https://oeis.org/A084613}; \url{https://oeis.org/A084615}.

\bibitem{hanna-triangles}
P.~D. Hanna, \emph{OEIS} A084600, A084602, A084604, A084606,
A084608, A084610, A084612 and A084614, definitions and formulas;
row formulas contributed by G.~C. Greubel, March 2023, and
G.~W. Adamson, May 2008.
\url{https://oeis.org/A084610}; \url{https://oeis.org/A084612};
\url{https://oeis.org/A084614}; \url{https://oeis.org/A084608}.

\bibitem{lean}
\emph{A195417RootLimit.lean},\newline declaration
\texttt{D5.S3.Analytic.Polynomial.A195417RootLimit.result},
revision \texttt{6bdbb68454659fcccffc7607b5fc292b2ff2fcdf}.
\href{https://github.com/the-omega-institute/trureturing/blob/6bdbb68454659fcccffc7607b5fc292b2ff2fcdf/D5/S3/Analytic/Polynomial/A195417RootLimit.lean}{Pinned Lean source}.
\end{thebibliography}
\bigskip
\noindent\small\textcopyright\ 2026 Haobo Ma and Wenlin Zhang.
Licensed under Creative Commons Attribution 4.0 International
(\href{https://creativecommons.org/licenses/by/4.0/}{CC BY 4.0}).
\end{document}
